\documentclass[11pt,a4paper]{article}
\usepackage[utf8]{inputenc}
\usepackage[T1]{fontenc}
\usepackage[margin=0.8in]{geometry}
\usepackage{amsmath,amssymb,booktabs,array,graphicx}
\usepackage{microtype,xcolor,placeins,float}
\usepackage[hidelinks]{hyperref}
\usepackage{caption}
\captionsetup[figure]{position=top,labelfont=bf,font=small,skip=4pt}
\captionsetup[table]{position=top,labelfont=bf,font=small,skip=4pt}
\setlength{\parindent}{1.15em}
\setlength{\parskip}{0.10em}
\setlength{\textfloatsep}{7pt plus 2pt minus 2pt}
\setlength{\floatsep}{6pt plus 2pt minus 2pt}
\setlength{\abovedisplayskip}{6pt}
\setlength{\belowdisplayskip}{6pt}
\setlength{\emergencystretch}{2em}
\newcommand{\CapitalSourceHash}{\texttt{3a76c95396024d98\ldots}}
\newcommand{\LoessDates}{1933, 2017}
\newcommand{\ParametricCandidateDates}{1945, 1971, 1989}
\newcommand{\IndependentKPUdmax}{50.118}
\newcommand{\IndependentKPCritical}{12.25}
\newcommand{\DynamicBreakDate}{2005}
\newcommand{\DynamicBreakCI}{1997--2019}
\newcommand{\DynamicUDmax}{34.064}
\newcommand{\DynamicUDmaxCV}{14.23}
\newcommand{\DynamicLag}{1}
\newcommand{\LongRunVerdict}{not identified}
\newcommand{\AccumulationVerdict}{not identified}


\title{\textbf{Structural Change without an Identified
Capital-Composition Elasticity: Chile, 1920--2024}}
\author{Chapter 2 Empirical Report}
\date{July 2026}

\begin{document}
\maketitle

\begin{abstract}
This paper asks whether Chile's capital stocks identify changes in the
long-run elasticity of nonresidential structures with respect to machinery
and equipment. The canonical annual panel contains 105 observations from
1920 through 2024. The answer is definitive but negative. Neither stable nor
break-adjusted cointegration is established, so the data identify no
\(\beta_{\mathrm{ME}}\) and no period-specific
\(\beta_{\mathrm{ME},j}\). A separate model of annual accumulation detects a
change in 2005, with an analytic 95-percent interval of 1997--2019, but the
growth rates do not meet the declared piecewise-stationarity condition.
Consequently, the impact and three-year accumulation responses
\(\eta_{j,0}\) and \(\eta_{j,3}\) are also not identified. The capital data
therefore reveal historical changes in composition and in the conditional
growth relation, but they do not isolate a long-run or short-run elasticity.
This distinction prevents persistent stock trends, fitted regression breaks,
and descriptive curvature from being mistaken for changes in technique.
\end{abstract}

\FloatBarrier
\section{The answer and its economic meaning}

Chile's machinery and equipment stock \(K^{\mathrm{ME}}\) and its
nonresidential structures stock \(K^{\mathrm{NR}}\) represent different
dimensions of accumulation. Machinery is treated here as the intensive
margin through which production is reorganized, while structures form the
plant envelope within which those techniques operate. This is a maintained
analytical restriction, not a claim that the capital-stock data directly
measure technology. The empirical question is whether the two observed
stocks contain enough independent variation to measure their relationship.

Their log difference defines capital composition:
\begin{equation}
  \kappa_t=k_t^{\mathrm{ME}}-k_t^{\mathrm{NR}}.
  \label{eq:kappa}
\end{equation}
Because this is an identity, the vector
\((k_t^{\mathrm{NR}},k_t^{\mathrm{ME}},\kappa_t)'\) is exactly singular.
The system must therefore be bivariate,
\(X_t=(k_t^{\mathrm{NR}},k_t^{\mathrm{ME}})'\), with \(\kappa_t\) used as
a derived contrast rather than as a third endogenous variable.

The long-run relation is normalized as
\begin{equation}
 k_t^{\mathrm{NR}}
 =\alpha_j+\beta_{\mathrm{ME},j}k_t^{\mathrm{ME}}+u_t,
 \qquad t\in(T_{j-1},T_j].
 \label{eq:long-run}
\end{equation}
Within a statistically admissible period \(j\), a 10-percent permanent
increase in machinery would correspond to approximately
\(10\beta_{\mathrm{ME},j}\) percent more structures. A coefficient below one
would mean that machinery outpaces structures; a coefficient above one would
mean that structures increase more than proportionally. That interpretation
requires a stationary long-run error \(u_t\). Without cointegration, a fitted
slope between two persistent stocks is not a long-run elasticity.

\begin{table}[h!]
\centering\small
\caption{Definitive empirical answer}
\label{tab:definitive-answer}
\begin{tabular}{p{2.4cm}p{2.5cm}p{2.7cm}p{4.9cm}}
\toprule
\textbf{Quantity} & \textbf{Status} & \textbf{Result} & \textbf{Economic meaning}\\
\midrule
Long-run $\beta_{\mathrm{ME}}$ & Not identified & not identified & the data do not identify a stable long-run structures elasticity with respect to machinery \\
Changes in $\beta_{\mathrm{ME},j}$ & Not identified & not identified & no structural period supports a period-specific cointegrating elasticity \\
Conditional accumulation $\eta_j$ & Not identified & not identified & growth-rate persistence prevents an elasticity interpretation \\
Dynamic-relation break & Diagnostic change & 2005 (95\% interval 1997--2019) & the conditional growth relation visibly changes, but its coefficients are not admissible elasticities \\
LOESS curvature dates & Descriptive only & descriptive inflection markers only & changes in the tendency of capital composition, not parametric breaks \\
\bottomrule
\multicolumn{4}{p{13.1cm}}{\footnotesize \textbf{Note:} A dated change in an observed relation is not automatically an elasticity change. Long-run elasticity requires cointegration and rank one; the conditional accumulation elasticity additionally requires stationary growth rates.}
\end{tabular}
\end{table}


Table~\ref{tab:definitive-answer} separates three findings that would
otherwise be easy to conflate. First, the capital stocks do not establish a
stable or period-specific long-run elasticity. Second, the annual growth
relation changes around 2005, but its coefficients do not pass the
stationarity condition required for an accumulation elasticity. Third, the
curvature of \(\kappa_t\) supplies descriptive historical markers only. The
paper therefore dates changes in observed capital behavior without assigning
an elasticity magnitude that the data cannot sustain.

\FloatBarrier
\section{Capital stocks and capital composition}

The source is the canonical Chile capital panel in constant 2003 CLP. All 105
annual observations from 1920 through 2024 reproduce, and the frozen source
hash begins \CapitalSourceHash. Older results based on 91 observations or on
a 1920--2010 endpoint are not evidence for this sample. Earlier files that
labeled ordinary least squares with Newey--West errors or cumulative
regressions as FM--OLS or IM--OLS are also excluded. No central result below
uses those coefficients.

\begin{figure}[h!]
\centering
\caption{Chile's machinery and nonresidential structures stocks, 1920--2024}
\label{fig:stocks}
\includegraphics[width=0.86\textwidth]{../../output/CL/SBVECM/S20_DUAL_ROUTE_BREAK_AUDIT/figures/fig1_capital_log_stocks.pdf}
\vspace{1pt}\par{\footnotesize \textbf{Note:} The figure plots the observed log
capital stocks before deterministic adjustment. Their common upward movement
motivates a long-run test but does not establish cointegration. Both stocks
cumulate investment and retirements, so persistence is partly built into their
measurement.}
\end{figure}

The common rise in Figure~\ref{fig:stocks} is economically important but
statistically insufficient. Capital stocks are accumulated magnitudes; each
observation embeds the preceding investment history. A high correlation can
therefore emerge even when the residual combination in
equation~\eqref{eq:long-run} remains nonstationary. This is why the empirical
burden lies first on integration order and cointegration, not on the size of a
levels-regression slope.

\begin{figure}[h!]
\centering
\caption{Observed capital composition and its LOESS tendency}
\label{fig:loess}
\includegraphics[width=0.86\textwidth]{../../output/CL/SBVECM/S20_DUAL_ROUTE_BREAK_AUDIT/figures/fig2_kappa_loess_route.pdf}
\vspace{1pt}\par{\footnotesize \textbf{Note:} Points are the observed
\(\kappa_t\); the continuous line is a LOESS fit. Zero-crossings of the
fit's second derivative occur near \LoessDates. These dates describe changes
in curvature. The fitted line and its derivatives are neither economic time
series nor inputs to the parametric break tests.}
\end{figure}

Figure~\ref{fig:loess} brings the historical movement down to the observed
composition ratio. The scatter shows the data; the line summarizes its local
tendency. The second derivative is used only to mark changes in curvature.
It is never unit-root tested, never treated as a stochastic series, and never
used to choose the formal coefficient periods. Agreement or disagreement
between a curvature date and a parametric date is descriptive comparison,
not mutual validation.

\FloatBarrier
\section{What the structural-break tests establish}

A structural break can concern the path of one stock, the deterministic
component of a relation, the slope of a fitted regression, or the
cointegrating vector itself. These are different empirical objects.
Table~\ref{tab:method-map} states what each procedure can and cannot support.

\begin{table}[h!]
\centering\scriptsize
\caption{What the structural-break and cointegration procedures establish}
\label{tab:method-map}
\begin{tabular}{p{3.3cm}p{4.6cm}p{5.0cm}}
\toprule
\textbf{Procedure} & \textbf{Economic-statistical question} & \textbf{Permitted claim}\\
\midrule
Bai--Perron & Where do coefficients in a fitted regression change? & Candidate regression dates; not cointegration \\
Zivot--Andrews & Does one series reject a unit root after an unknown break? & Integration-order evidence for one series \\
Gregory--Hansen & Does cointegration appear after one unknown regime change? & One-break cointegration evidence \\
Kejriwal--Perron & Are coefficients stable in an already cointegrated regression? & Coefficient stability conditional on cointegration \\
Johansen--Mosconi--Nielsen & What is the rank with known deterministic breaks? & Conditional rank and vector restrictions \\
L"utkepohl--Saikkonen--Trenkler & What is the rank after one unknown system shift? & Conditional rank with one unknown level shift \\
Hansen--Johansen & Is an established cointegrating vector recursively constant? & Parameter constancy after rank is established \\
Phillips--Ouliaris & Does a residual combination reject no cointegration? & Stable cointegration evidence \\
Lee--Strazicich & Does one series reject a unit root with two admissible breaks? & Break-aware integration-order evidence \\
\texttt{mbreaks} & How are Bai--Perron statistics and intervals computed? & Implementation, not a separate economic test \\
\bottomrule
\multicolumn{3}{p{13.1cm}}{\footnotesize \textbf{Note:} Each procedure concerns a different statistical object. Evidence about a single series, a fitted regression, a cointegrating residual, and system rank cannot be substituted for one another.}
\end{tabular}
\end{table}


\subsection{From an individual series to a long-run relation}

Zivot--Andrews asks whether a single series rejects a unit root once one
unknown break is admitted under the alternative. In a stylized intercept-and-
trend specification,
\begin{equation}
 y_t=\mu+\theta DU_t(\lambda)+\beta t+\gamma DT_t(\lambda)
      +\rho y_{t-1}+\sum_{i=1}^{p}c_i\Delta y_{t-i}+e_t ,
 \label{eq:za}
\end{equation}
the break fraction \(\lambda\) is selected where the unit-root statistic is
most adverse to the null. Rejection informs the integration classification
of \(y_t\); it does not show that the slope linking
\(k_t^{\mathrm{NR}}\) to \(k_t^{\mathrm{ME}}\) changed.
Lee--Strazicich provides a complementary two-break LM test that permits
breaks under both null and alternative. Neither test, by itself, identifies
\(\beta_{\mathrm{ME},j}\).

Bai--Perron asks whether coefficients in a stable regression are constant
across unknown segments. For \(m\) breaks,
\begin{equation}
 y_t=x_t'\vartheta+z_t'\delta_j+e_t,
 \qquad t=T_{j-1}+1,\ldots,T_j ,
 \label{eq:bp}
\end{equation}
it minimizes the global residual sum of squares subject to a minimum segment
length and tests zero against one or more breaks. The procedure is highly
useful for locating candidate changes and constructing confidence intervals.
In a levels regression with integrated variables, however, ordinary
Bai--Perron inference does not establish that the residual is stationary.
Its dates cannot, on their own, divide the long-run elasticity.

\begin{figure}[h!]
\centering
\caption{Bai--Perron segmented fits: candidate changes in the capital relation}
\label{fig:bp}
\includegraphics[width=0.92\textwidth]{../../output/CL/SBVECM/S20_DUAL_ROUTE_BREAK_AUDIT/figures/fig3_bai_perron_segmented_fits.pdf}
\vspace{1pt}\par{\footnotesize \textbf{Note:} Segmented fits and confidence
intervals reproduce the candidate evidence for the stock equations,
\(\kappa_t\), and the normalized capital relation. The principal parametric
candidate dates are \ParametricCandidateDates. They locate instability in
fitted relations; they do not establish cointegration or a change in
\(\beta_{\mathrm{ME}}\).}
\end{figure}

Figure~\ref{fig:bp} restores the break graph omitted from the earlier report.
Its economic value is to show that a constant descriptive relation is
implausible over the century. Its inferential limit is equally important:
smooth persistent stocks can generate sharply segmented fits even when no
stationary long-run error exists.

\subsection{From candidate instability to cointegration}

Gregory--Hansen tests the null of no cointegration against cointegration with
one unknown level, trend, or regime change. A regime-shift version can be
written
\begin{equation}
 k_t^{\mathrm{NR}}=\alpha_1+\alpha_2DU_t(\lambda)
 +\beta_1 k_t^{\mathrm{ME}}
 +\beta_2DU_t(\lambda)k_t^{\mathrm{ME}}+u_t .
 \label{eq:gh}
\end{equation}
The statistic takes the most favorable residual-based unit-root value over
the trimmed set of possible dates. Rejection would establish one-break
cointegration under the chosen deterministic form; the estimated date would
then be eligible for system analysis.

Kejriwal--Perron addresses multiple coefficient changes in an already
cointegrated regression with integrated regressors. It therefore asks the
right stability question only after cointegration is established. The
UDmax statistic tests stability against an unknown number of breaks, and
sequential supF tests select their count subject to the trimming rule. A
large statistic can reject coefficient stability conditional on
cointegration; it cannot supply the missing cointegration prerequisite.

Johansen--Mosconi--Nielsen then provides rank inference with known
deterministic breaks, while L\"utkepohl--Saikkonen--Trenkler treats one
unknown system-level shift. Hansen--Johansen recursive statistics concern
the constancy of a cointegrating vector after its rank is established.
These system procedures answer whether equation~\eqref{eq:long-run} belongs
to a coherent adjustment system; they are not substitutes for the
single-series or residual-cointegration evidence.

\FloatBarrier
\section{Long-run cointegration and
\texorpdfstring{\(\beta_{\mathrm{ME}}\)}{beta ME}}

The full integration audit applies ADF, Phillips--Perron, KPSS, ERS--GLS,
Schmidt--Phillips, and Zivot--Andrews tests to
\(k_t^{\mathrm{NR}}\), \(k_t^{\mathrm{ME}}\), and \(\kappa_t\) in levels,
first differences, and second differences, supplemented by break-aware
tests. The primary deterministic cases do not produce a common, unambiguous
\(I(1)\) classification for the two stocks. Conflicting evidence is retained
as inconclusive rather than forced into an \(I(1)\) designation.

\begin{table}[h!]
\centering\small
\caption{Long-run identification evidence}
\label{tab:long-run}
\begin{tabular}{p{3.0cm}p{3.25cm}p{1.5cm}p{5.2cm}}
\toprule
\textbf{Evidence} & \textbf{Result} & \textbf{Gate} & \textbf{Permitted conclusion}\\
\midrule
Raw integration order & machinery stock inconclusive; structures stock inconclusive & No & raw capital orders remain inconclusive \\
Gregory--Hansen & minimum -4.778 against -4.95 & No & one-break cointegration is not supported \\
Phillips--Ouliaris & 0 of 6 reject & No & stable cointegration is not supported \\
No-break Johansen & rank=2 & No & the system does not supply an admissible rank-one relation \\
Independent Kejriwal--Perron & UDmax 50.118 against 12.25 & No & stability statistic cannot define coefficient periods \\
Validated long-run breaks & none & No & no date defines a long-run coefficient period \\
Stable $\beta_{\mathrm{ME}}$ & not identified & No & not run: integration, cointegration, and rank prerequisites fail \\
Period $\beta_{\mathrm{ME},j}$ & not identified & No & not run: integration, cointegration, and rank prerequisites fail \\
\bottomrule
\multicolumn{4}{p{13.1cm}}{\footnotesize \textbf{Note:} The independent Kejriwal--Perron statistic exceeds its critical value, but the procedure assumes cointegration. Since neither Gregory--Hansen nor Phillips--Ouliaris establishes cointegration, the stability rejection cannot define coefficient periods.}
\end{tabular}
\end{table}


The residual and system evidence in Table~\ref{tab:long-run} resolves the
question. The best Gregory--Hansen statistic is \(-4.778\), which does not
cross its relevant 5-percent critical value of \(-4.95\). None of the six
Phillips--Ouliaris variants rejects no cointegration. The no-break Johansen
system selects rank two rather than the single cointegrating vector required
to identify equation~\eqref{eq:long-run}.

The independent \texttt{mbreaks} implementation produces
\(\mathrm{UDmax}=\IndependentKPUdmax\), above the 5-percent critical value
\(\IndependentKPCritical\). This corrects the earlier custom computation,
which understated the instability evidence. Yet the correction does not
create coefficient periods: the Kejriwal--Perron distribution presumes a
cointegrated regression, while the preceding tests do not establish one.
The statistic therefore says that a stable cointegrating regression would
be rejected if such a regression were admissible; it cannot prove that the
capital stocks cointegrate.

No date consequently passes the full sequence of integration, break-aware
cointegration, rank-one, and residual-dynamics conditions. The actual
\texttt{cointRegIM} estimator is intentionally not run, because a
super-consistent coefficient is not meaningful after those prerequisites
fail. FM--OLS and DOLS cannot repair the failure and remain excluded from the
central evidence. There is therefore no admissible
\(\beta_{\mathrm{ME}}\), no \(\beta_{\mathrm{ME},j}\), and no statistical
basis for testing whether a reported period coefficient equals one.

The same logic governs \(\kappa_t\). Global stationarity would require the
unit-vector restriction \(\beta_{\mathrm{ME}}=1\), no slope changes, and no
permanent deterministic shift in \(\kappa_t\). Piecewise stationarity would
permit level shifts but still require every regime-specific machinery
coefficient to equal one. Since no cointegrating vector is identified, these
restrictions cannot be evaluated as long-run restrictions. Visual proximity
of a fitted coefficient to one would not establish stationarity.

\FloatBarrier
\section{Stationary accumulation and
\texorpdfstring{\(\eta_j\)}{eta j}}

Long-run non-identification does not rule out a separate relation between
annual accumulation rates. The relevant stationary model is
\begin{equation}
\Delta k_t^{\mathrm{NR}}
=c_j+\sum_{\ell=1}^{p}a_\ell\Delta k_{t-\ell}^{\mathrm{NR}}
 +\sum_{\ell=0}^{p}b_{j\ell}\Delta k_{t-\ell}^{\mathrm{ME}}+e_t .
\label{eq:dynamic}
\end{equation}
Here \(\eta_{j,0}=b_{j0}\) would measure the conditional impact response of
structures growth to machinery growth, while
\(\eta_{j,3}\) would cumulate the dynamic response through three years.
These are conditional accumulation elasticities, not the long-run
cointegrating coefficient in equation~\eqref{eq:long-run}.

The lag order is chosen on a common sample. BIC selects \(p=\DynamicLag\);
the resulting equation passes the four-lag residual serial-correlation test
with \(p=0.504\). Allowing the intercept and machinery-response coefficients
to change, the break test gives
\(\mathrm{UDmax}=\DynamicUDmax\), above its 5-percent critical value
\(\DynamicUDmaxCV\). The sequential procedure rejects zero against one break
but not one against two, locating the change in \DynamicBreakDate\ with a
95-percent interval of \DynamicBreakCI.

\begin{table}[h!]
\centering\small
\caption{Break tests for the dynamic accumulation relation}
\label{tab:dynamic-break}
\begin{tabular}{p{2.3cm}p{1.7cm}p{1.7cm}p{1.4cm}p{5.5cm}}
\toprule
\textbf{Test} & \textbf{Statistic} & \textbf{5\% critical} & \textbf{Reject} & \textbf{Sequential implication}\\
\midrule
UDmax & 34.064 & 14.23 & Yes & 2005 (1997--2019) \\
supF(1|0) & 34.064 & 13.98 & Yes & continue \\
supF(2|1) & 9.503 & 15.72 & No & stop \\
supF(3|2) & 15.324 & 16.83 & No & stop \\
\bottomrule
\multicolumn{5}{p{13.1cm}}{\footnotesize \textbf{Note:} The dynamic equation uses one annual lag, selected by BIC on a common sample, and passes a four-lag residual serial-correlation test. These statistics locate a visible conditional change. They do not become elasticity evidence because the growth-rate stationarity condition fails.}
\end{tabular}
\end{table}


This is strong evidence that the conditional growth relation is not constant.
It is not yet an elasticity result. The statistical meaning of the
coefficients in equation~\eqref{eq:dynamic} depends on stationary growth
rates within the deterministic specification used to define the break.
The complete battery is therefore repeated after allowing a 2005 mean shift.

\begin{table}[h!]
\centering\scriptsize
\caption{Stationarity evidence for the capital-growth rates after the 2005 mean shift}
\label{tab:growth-stationarity}
\begin{tabular}{p{1.35cm}p{1.35cm}p{1.35cm}p{1.35cm}p{1.35cm}p{1.35cm}p{1.35cm}p{1.75cm}}
\toprule
\textbf{Series} & \textbf{ADF} & \textbf{PP} & \textbf{KPSS} & \textbf{ERS} & \textbf{SP} & \textbf{ZA} & \textbf{Conclusion}\\
\midrule
$\Delta k^{\mathrm{NR}}$ & stationary & stationary & stationary & conflict & unit root & unit root & Inconclusive \\
$\Delta k^{\mathrm{ME}}$ & stationary & conflict & stationary & unit root & unit root & unit root & Inconclusive \\
\bottomrule
\multicolumn{8}{p{13.1cm}}{\footnotesize \textbf{Note:} A family verdict summarizes its declared primary deterministic cases so that densely parameterized test families receive no extra weight. Four concordant families are required. The disagreement therefore remains substantive rather than being resolved by majority choice among individual configurations.}
\end{tabular}
\end{table}


For \(\Delta k_t^{\mathrm{NR}}\), three of the six primary test families
support stationarity; ERS is conflicting, while Schmidt--Phillips and
Zivot--Andrews retain a unit root. For
\(\Delta k_t^{\mathrm{ME}}\), only ADF and KPSS support stationarity;
Phillips--Perron conflicts and the remaining families retain a unit root.
The declared requirement of four concordant families is not met for either
series, and the two-break Lee--Strazicich tests do not reverse that verdict.
The evidence remains inconclusive instead of being resolved by counting the
many specifications inside one test family.

\begin{figure}[h!]
\centering
\caption{Annual capital accumulation and the 2005 diagnostic break}
\label{fig:accumulation}
\includegraphics[width=0.90\textwidth]{../../output/CL/SBVECM/S35_DYNAMIC_ACCUMULATION_RESPONSE/figures/fig8_accumulation_evidence.pdf}
\vspace{1pt}\par{\footnotesize \textbf{Note:} The left panel plots observed
annual log changes; the right panel displays their contemporaneous
association. The vertical line marks the diagnostic 2005 change. No fitted
response line is shown because the stationarity condition needed to interpret
its slope as \(\eta_j\) is not satisfied.}
\end{figure}

\begin{table}[h!]
\centering\small
\caption{Can the conditional accumulation elasticity be reported?}
\label{tab:eta}
\begin{tabular}{p{2.0cm}p{0.8cm}p{1.5cm}p{1.8cm}p{2.0cm}p{4.7cm}}
\toprule
\textbf{Span} & \textbf{$N$} & \textbf{$\eta_{j,0}$} & \textbf{$\eta_{j,3}$} & \textbf{Adjacent test} & \textbf{Empirical verdict}\\
\midrule
1922-2005 & 84 & -- & -- & -- & Not identified: growth-rate stationarity is inconclusive \\
2006-2024 & 19 & -- & -- & -- & Not identified: growth-rate stationarity is inconclusive \\
\bottomrule
\multicolumn{6}{p{13.1cm}}{\footnotesize \textbf{Note:} $\eta_{j,0}$ would be the impact response and $\eta_{j,3}$ the cumulative three-year response. Both are intentionally blank because the growth rates do not pass the declared piecewise-stationarity rule. The 2005 break remains a feature of the observed dynamic relation, not an estimated elasticity change.}
\end{tabular}
\end{table}


The blank coefficient cells in Table~\ref{tab:eta} are empirical results,
not missing calculations. Estimating separate period regressions would
manufacture magnitudes after the governing condition failed, especially in
the 19-observation post-2005 period. The conditional bivariate VAR and
Granger-exclusion tests are also not reported because they require the same
stationary system. Resampling cannot remedy this problem: bootstrap inference
can approximate the distribution of an admissible statistic, but it cannot
turn persistent growth rates into stationary variables. No bootstrap is
therefore necessary or appropriate for the reported verdict.

\FloatBarrier
\section{Historical interpretation}

The evidence establishes structural change in Chilean capital accumulation
without identifying its elasticity. The capital stocks visibly change their
relative tendency over the century, ordinary segmented regressions reject
constancy, and the dynamic growth relation changes around 2005. These
findings reject a timeless descriptive mapping from machinery into
structures. They do not reveal by how much a permanent machinery expansion
changes structures in the long run, nor do they supply a stationary
short-run response.

This boundary matters for the choice-of-technique interpretation. A rising
\(\kappa_t\) is consistent with machinery deepening relative to the plant
envelope, but it does not by itself identify the behavioral mechanism.
Capital-stock construction combines investment, retirement, relative
prices, and inherited stocks. With only the two stock series, secular
measurement persistence and economic adjustment cannot be separated tightly
enough to assign a coefficient magnitude. The maintained intensive/extensive
distinction supplies the economic question; it cannot supply the missing
statistical variation.

The 2005 change has a similarly bounded interpretation. It indicates that
the observed conditional association between the two accumulation rates
changes within a broad 1997--2019 interval. The interval is too wide for a
single institutional date to be inferred from the capital data alone, and
the nonstationary evidence blocks a parameter interpretation. The result is
historically suggestive but neither causal nor an estimate of technique.

\FloatBarrier
\section{Conclusion}

The canonical 1920--2024 Chile capital panel does not identify the order of
magnitude of \(\beta_{\mathrm{ME}}\), either as one stable coefficient or as
period-specific coefficients across structural breaks. The decisive reason
is not the absence of visible breaks. On the contrary, instability is strong:
the independent Kejriwal--Perron statistic and the segmented fits reject
constancy. The reason is that stable and one-break cointegration are not
established and the system does not supply an admissible rank-one vector.
Candidate dates therefore cannot become long-run coefficient periods.

The separate accumulation analysis reaches the other predetermined outcome.
A single dynamic change is detected in 2005, but neither capital-growth rate
passes the declared piecewise-stationarity condition. The evidence therefore
does not identify \(\eta_{j,0}\) or \(\eta_{j,3}\). Chile's capital data
establish historical change in composition and accumulation, but they
identify neither a long-run elasticity nor a conditional accumulation
elasticity. That non-identification is the definitive answer supplied by the
available capital stocks.

\FloatBarrier
\appendix
\section{Compact methodological record}

\subsection{Integration and break adjudication}

The unit-root audit records deterministic case, lag, statistic, and critical
value for ADF, Phillips--Perron, KPSS, ERS--GLS,
Schmidt--Phillips, and Zivot--Andrews tests in levels, first differences, and
second differences. Test families receive one primary verdict per
deterministic case so that families with more tuning choices do not dominate.
Break-aware Lee--Strazicich results supplement, rather than replace, this
evidence. A deterministic step may absorb a level shift; it is not presumed
to convert a genuine \(I(2)\) process into \(I(1)\).

The parametric evidence uses a 15-percent minimum segment and no more than
three breaks. Ordinary Bai--Perron dates are candidate changes.
Gregory--Hansen supplies one-break residual cointegration evidence.
Kejriwal--Perron supplies multiple-break stability evidence conditional on
cointegration. Known deterministic shifts belong in the
Johansen--Mosconi--Nielsen system, while an unknown system shift belongs in
the L\"utkepohl--Saikkonen--Trenkler procedure. No LOESS curvature date enters
these computations.

\subsection{Coefficient construction and inference}

Had validated long-run breaks existed, period coefficients would have been
estimated jointly with initial-break-state-centered interactions:
\begin{equation}
k_t^{\mathrm{NR}}=\alpha_1+\beta_1 k_t^{\mathrm{ME}}
+\sum_{b=1}^{m}\gamma_bD_{b,t}
+\sum_{b=1}^{m}\theta_bD_{b,t}
  \bigl(k_t^{\mathrm{ME}}-k_{T_b}^{\mathrm{ME}}\bigr)+u_t ,
\end{equation}
where \(\beta_{\mathrm{ME},1}=\beta_1\) and
\(\beta_{\mathrm{ME},j}=\beta_1+\sum_{b<j}\theta_b\).
Here \(\gamma_b\) denotes deterministic step controls in the single equation;
the system form uses \(\boldsymbol{\Gamma}\mathbf D_t\).
Uncentered and sample-mean-centered versions would be numerical
comparators assessed by confidence-band overlap and condition numbers.
Residualization is not used. Since the cointegration conditions fail, these
coefficients and their cumulative covariance are not estimated.

Analytic break-aware and HAC inference is the default. Restricted-null
residual bootstrap inference would be reserved for an eligible
estimated-break system lacking analytic rank inference, material
break-selection uncertainty not covered analytically, or an indispensable
period with at most 30 observations. None of those situations arises after
the identification conditions fail, so the number of bootstrap replications
used is zero.

\subsection{System interpretation}

The bivariate error-correction representation underlying the system checks is
\begin{equation}
\Delta X_t=\boldsymbol{\Pi}X_{t-1}
+\sum_{i=1}^{p-1}\boldsymbol{\Gamma}_i\Delta X_{t-i}
+\boldsymbol{\Gamma}\mathbf D_t+\varepsilon_t,
\qquad
\boldsymbol{\Pi}=\boldsymbol{\alpha}\boldsymbol{\beta}' .
\end{equation}
For two persistent capital stocks, rank zero means that no stationary
long-run combination enters the system. Rank one is the economically relevant
case: \(\boldsymbol{\beta}\) contains the normalized capital-composition
relation and \(\boldsymbol{\alpha}\) records how each stock responds to its
disequilibrium. Rank two makes the levels system full rank and therefore does
not deliver the single common-trend relation required to interpret
\(\beta_{\mathrm{ME}}\). The no-break rank-two result is consequently not
evidence of a stronger form of the desired cointegration; it shows that the
proposed \(I(1)\), rank-one representation is not supported.

Known step shifts enter through \(\boldsymbol{\Gamma}\mathbf D_t\), while
slope changes belong in the cointegrating vector and require break-aware rank
inference. A deterministic shift can alter the mean of \(\kappa_t\), but it
cannot be used mechanically to remove stochastic persistence. Adjustment
coefficients, weak-exogeneity tests, recursive constancy, and companion-root
diagnostics become interpretable only after an admissible rank is established.
They are therefore not reported as isolated system statistics after the
rank-one condition fails.

\FloatBarrier
\clearpage
\section{References}

\begin{description}
\setlength{\itemsep}{2pt}
\item Bai, J., and P. Perron. 2003. ``Computation and Analysis of Multiple
Structural Change Models.'' \emph{Journal of Applied Econometrics}
18(1): 1--22.
\item Gregory, A. W., and B. E. Hansen. 1996. ``Residual-Based Tests for
Cointegration in Models with Regime Shifts.'' \emph{Journal of Econometrics}
70(1): 99--126.
\item Hansen, H., and S. Johansen. 1999. ``Some Tests for Parameter Constancy
in Cointegrated VAR-Models.'' \emph{Econometrics Journal} 2(2): 306--333.
\item Johansen, S., R. Mosconi, and B. Nielsen. 2000. ``Cointegration Analysis
in the Presence of Structural Breaks in the Deterministic Trend.''
\emph{Econometrics Journal} 3(2): 216--249.
\item Kejriwal, M., and P. Perron. 2010. ``Testing for Multiple Structural
Changes in Cointegrated Regression Models.'' \emph{Journal of Business and
Economic Statistics} 28(4): 503--522.
\item Lee, J., and M. C. Strazicich. 2003. ``Minimum Lagrange Multiplier Unit
Root Test with Two Structural Breaks.'' \emph{Review of Economics and
Statistics} 85(4): 1082--1089.
\item L\"utkepohl, H., P. Saikkonen, and C. Trenkler. 2004. ``Testing for the
Cointegrating Rank of a VAR Process with Level Shift at Unknown Time.''
\emph{Econometrica} 72(2): 647--662.
\item Nguyen, T. H., Y. Yamamoto, and P. Perron. 2025. \emph{mbreaks:
Multiple Break Testing and Estimation}. R package version 1.0.1.
\item Phillips, P. C. B., and S. Ouliaris. 1990. ``Asymptotic Properties of
Residual Based Tests for Cointegration.'' \emph{Econometrica}
58(1): 165--193.
\item Zivot, E., and D. W. K. Andrews. 1992. ``Further Evidence on the Great
Crash, the Oil-Price Shock, and the Unit-Root Hypothesis.''
\emph{Journal of Business and Economic Statistics} 10(3): 251--270.
\end{description}

\end{document}
